\cotefile{10}
\scanpage{77}
quotient de $\handwritten{\overline}{\handwritten{\underline}{Q}}^{\handwritten{*}}$ par le sous-groupe des unités, donc canoniquement isomorphe à $\handwritten{\underline}{Q}$, par l'isomorphisme
\[
\handwritten{\overline}{\handwritten{\underline}{Q}}^{\handwritten{*}}/\text{unités}\handwritten{\longrightarrow}\handwritten{\underline}{Q}
\]
donnés par la valuation p-adique $v_p$, normalisée par la condition
\[
v_p(p)=1\quad .
\]
\unsure{A|L}'opération de $\deleted{\unsure{\pi_1(\underline{k}|\pi_1(\underline{K}|\pi(\underline{k}}}$ $\operatorname{Gal}(\handwritten{\overline}{\handwritten{\underline}{Q}}_p/\handwritten{\underline}{Q}_p)$ sur ce groupe quotient $Q$ est \unsure{triviale|trivialé|trivialè}, de sorte que $H$ est \handwritten{\uline}{diagonalisable} de groupe des caractères $\handwritten{\underline}{Q}$ :
\[
H=D(\handwritten{\underline}{Q})\quad .
\]
L'augmentation de Tate de $G$ correspond à l'homomorphisme
\[
\handwritten{\underline}{Z}\handwritten{\longrightarrow}\handwritten{\overline}{\handwritten{\underline}{Q}}^{\handwritten{*}},\quad 1\handwritten{\longmapsto}p\quad ,
\]
et l'augmentation induite sur $H$ correspond à l'\handwritten{\uline}{inclusion} canonique
\[
\handwritten{\underline}{Z}\handwritten{\lhook\joinrel\longrightarrow}\handwritten{\underline}{Q}
\]
(\unsure{grâee|grâce} à notre choix de la valuation \handwritten{\uline}{normalisée} comme indiquée).

\unsure{B|8}.2. Il faut maintenant expliciter l'homomorphisme de gerbes
\[
\handwritten{i:}\ \handwritten{\underline}{G}(\unsure{\handwritten{\underline}{F}|\handwritten{\underline}{\mathbb{F}}}_p)\handwritten{\longrightarrow}\handwritten{\underline{\Pi}}_p(\handwritten{\underline}{F}_p)
\tag*{\handwritten{(8.2.1)}}
\]
correspondant à l'inclusion
\[
i:\quad \handwritten{\underline}{Q}_p\text{-faisceaux constructibles}\handwritten{\lhook\joinrel\longrightarrow}\operatorname{Fcriso}(\handwritten{\underline}{F}_p)
\]
de 5.2. On a vu que l'image essentielle est exactement égale à la sous-catégorie \inserted{pleine} $J(\handwritten{\underline}{Q}_p\text{-faisceaux constructibles})$, \deleted{qui est} de sorte que $i$ s'identifie à une\deleted{foncteur} $\handwritten{\otimes}$-équivalence de la catégorie des $\handwritten{\underline}{Q}_p$-faisceaux constructibles avec elle-même. On constate alors que l'homomorphisme correspondant sur les liens $\deleted{\unsure{H\ H|H\to H}}$ $G'\handwritten{\longrightarrow}G'$ est l'identité, \deleted{et que} de sorte que $i$ est défini par un torseur sous $G'$, et que ce torseur sous $G'$ peut se décrire ainsi (aurait pu passer à la fin de 7, dans le cas d'un $k$ quelconque ...). On a un homomorphisme canonique
\[
\handwritten{\pi}_1(\handwritten{\underline}{F}_p^{\deleted{\handwritten{\unsure{*|\times}}}})=\handwritten{\hat}{\handwritten{\underline}{Z}}^{\deleted{\handwritten{\unsure{*|\times}}}}\handwritten{\longrightarrow}G'(\handwritten{\underline}{Q}_p),\quad \handwritten{\text{i.e. }(\pi_1(F_p))_{\mathbb{Q}_p}\longrightarrow G'}
\tag*{\handwritten{(8.2.2)}}
\]
d'où un \deleted{\unsure{homomorphi|homomorphix}} foncteur
\[
\text{torseurs}^{\deleted{\unsure{\underline{Q}_p|\underline{Q}_p\text{xxx}}}}\text{ sous }(\handwritten{\hat}{\handwritten{\underline}{Z}})_{\handwritten{\underline}{Q}_p}\handwritten{\longrightarrow}\text{torseurs sous }G'\ \deleted{\handwritten{\unsure{\xi|\varepsilon|,|.}}}\quad .
\]
% archive sheet number: 76
